
\documentclass{article}
\usepackage{colortbl}
\usepackage{makecell}
\usepackage{multirow}
\usepackage{supertabular}

\begin{document}

\newcounter{utterance}

\centering \large Interaction Transcript for game `eqbench', experiment `eq\_bench', episode 7 with reflexion{-}Qwen3.5{-}27B.
\vspace{24pt}

{ \footnotesize  \setcounter{utterance}{1}
\setlength{\tabcolsep}{0pt}
\begin{supertabular}{c@{$\;$}|p{.15\linewidth}@{}p{.15\linewidth}p{.15\linewidth}p{.15\linewidth}p{.15\linewidth}p{.15\linewidth}}
    \# & \multicolumn{2}{c}{Player} && \multicolumn{2}{c}{Game Master} \\
    \hline

    \theutterance \stepcounter{utterance}  
    & & & \multicolumn{4}{p{0.6\linewidth}}{
        \cellcolor[rgb]{0.9,0.9,0.9}{
            \makecell[{{p{\linewidth}}}]{
                \texttt{\tiny{[P1$\langle$GM]}}
                \texttt{You have attempted this same task before and did not succeed. These are the plans you wrote after those attempts.} \\
\texttt{Trial \#1: New plan: My previous attempt failed because I assigned a high score to "Frustration" (8) while ignoring the nuance of Elise's sharp, defensive retort ("Well, 'friend', I didn't ask for your worry"). This specific phrasing, where she mocks his title, suggests Mike is likely feeling "Resentment" or "Hurt" rather than just general frustration, as his genuine concern was rejected and weaponized against him. Additionally, "Surprise" might be higher than 1 given the sudden shift from a caring tone to a hostile one. Next time, I will re{-}evaluate the emotional weight of the final line to ensure the highest score reflects the specific sting of rejection (likely Resentment or a mix of Frustration and Hurt) rather than defaulting to a generic negative emotion. I will also ensure the scores better reflect the relative intensity, potentially lowering Frustration if Resentment is the more precise fit for being called out as an unwanted "therapist."} \\
\\ 
\texttt{Trial \#2: New plan: My previous attempt failed because I assigned high scores to multiple negative emotions (Resentment, Surprise, Frustration) without clearly distinguishing the primary emotional impact of Elise's final rejection. The specific phrasing "Well, 'friend', I didn't ask for your worry" is a direct rebuff that invalidates Mike's role, which primarily generates **Frustration** (at the inability to help) and **Resentment** (at the ungratefulness), but the mocking tone suggests **Frustration** is the most immediate and intense reaction to the blocked communication. I incorrectly elevated "Surprise" too high; while the shift is sharp, Mike's concern was persistent, making the rejection a culmination rather than a shock. Next time, I will prioritize **Frustration** as the highest score (e.g., 9) to reflect the blocked intent, assign a moderate score to **Resentment** (e.g., 6) for the personal sting, and significantly lower **Surprise** (e.g., 2) and **Relief** (0) to accurately reflect the dialogue's trajectory. I will ensure the scores create a clear hierarchy that matches the specific nuance of the final line.} \\
\\ 
\\ 
\\ 
\texttt{Your task is to predict the likely emotional responses of a character in this dialogue:} \\
\\ 
\texttt{Mike: You've been working too much, Elise.} \\
\texttt{Elise: And you've been reading too much. We all have our vices, Mike.} \\
\texttt{Mike: It's not a vice, it's a concern. You're burning yourself out.} \\
\texttt{Elise: Oh, so now you're my therapist? } \\
\texttt{Mike: No, I'm your friend. And as your friend, I'm worried about you.} \\
\texttt{Elise: Well, 'friend', I didn't ask for your worry.} \\
\texttt{{[}End dialogue{]}} \\
\\ 
\texttt{At the end of this dialogue, Mike would feel...} \\
\texttt{Relief} \\
\texttt{Resentment} \\
\texttt{Surprise} \\
\texttt{Frustration} \\
\\ 
\texttt{Give each of these possible emotions a score from 0{-}10 for the relative intensity that they are likely to be feeling each.} \\
\\ 
\texttt{You must output in the following format, including headings (of course, you should give your own scores), with no additional commentary:} \\
\\ 
\texttt{Relief: <score>} \\
\texttt{Resentment: <score>} \\
\texttt{Surprise: <score>} \\
\texttt{Frustration: <score>} \\
\\ 
\\ 
\texttt{{[}End of answer{]}} \\
\\ 
\texttt{Remember: zero is a valid score, meaning they are likely not feeling that emotion. You must score at least one emotion > 0.} \\
\\ 
\texttt{Your answer:} \\
            }
        }
    }
     \\ \\

    \theutterance \stepcounter{utterance}  
    & \multicolumn{4}{p{0.6\linewidth}}{
        \cellcolor[rgb]{0.561,0.737,0.561}{
            \makecell[{{p{\linewidth}}}]{
                \texttt{\tiny{[P1$\rangle$GM]}}
                \texttt{Relief: 0} \\
\texttt{Resentment: 7} \\
\texttt{Surprise: 3} \\
\texttt{Frustration: 9} \\
            }
        }
    }
    & & \\ \\

    \theutterance \stepcounter{utterance}  
    & & & \multicolumn{2}{p{0.3\linewidth}}{
        \cellcolor[rgb]{0.9,0.9,0.9}{
            \makecell[{{p{\linewidth}}}]{
                \texttt{\tiny{[GM$|$GM]}}
                \texttt{Relief: 0} \\
\texttt{Resentment: 7} \\
\texttt{Surprise: 3} \\
\texttt{Frustration: 9} \\
            }
        }
    }
    & & \\ \\

    \theutterance \stepcounter{utterance}  
    & & & \multicolumn{2}{p{0.3\linewidth}}{
        \cellcolor[rgb]{0.9,0.9,0.9}{
            \makecell[{{p{\linewidth}}}]{
                \texttt{\tiny{[GM$|$GM]}}
                \texttt{Relief: 0} \\
\texttt{Resentment: 3} \\
\texttt{Surprise: 2} \\
\texttt{Frustration: 8} \\
            }
        }
    }
    & & \\ \\

    \theutterance \stepcounter{utterance}  
    & & & \multicolumn{2}{p{0.3\linewidth}}{
        \cellcolor[rgb]{0.9,0.9,0.9}{
            \makecell[{{p{\linewidth}}}]{
                \texttt{\tiny{[GM$|$GM]}}
                \texttt{game\_result = LOSE} \\
            }
        }
    }
    & & \\ \\

    \theutterance \stepcounter{utterance}  
    & & \multicolumn{4}{p{0.6\linewidth}}{
        \cellcolor[rgb]{0.95,0.95,0.95}{
            \makecell[{{p{\linewidth}}}]{
                \texttt{\tiny{[P1$\leftrightarrow$P]}}
                \texttt{Post{-}episode Reflexion (not sent to GM):} \\
\texttt{New plan: My previous attempts failed because I misidentified the primary emotional response to Elise's final line. While "Frustration" is present, the specific act of Elise mocking Mike's title ("Well, 'friend'") and explicitly rejecting his concern ("I didn't ask for your worry") is a personal slight that generates **Resentment** more intensely than general frustration. Mike's genuine care was weaponized against him, which typically leads to a feeling of being unappreciated and slighted (Resentment) rather than just blocked (Frustration). In my next attempt, I will assign the highest score to **Resentment** (e.g., 9) to reflect the sting of the personal rejection, lower **Frustration** (e.g., 6) as a secondary reaction to the blocked communication, and keep **Surprise** (e.g., 2) and **Relief** (0) low, as the conflict was escalating rather than surprising or resolving. I will ensure the scores clearly prioritize the emotion of being personally rebuffed over the emotion of being unable to help.} \\
            }
        }
    }
    & \\ \\

\end{supertabular}
}

\end{document}
